\chapter{ブレーキ}
% full version: chapters/03_gaba.tex

数が2番目に多い種類のニューロン、\nt{GABA}をマシンに加えよう。その信号は隣をカチッの方へ押すのではなく、押し戻す。つまりブレーキである。数は多くない。大脳皮質では全ニューロンの5分の1から4分の1だ。だが、ブレーキなしではマシンは走れない。

ひとつ制約がある。「プラス」のニューロンは、自分で隣にブレーキをかけることはできない。負の結合が必要なら、仲介役を通して作るしかない。「プラス」が抑制性ニューロンを起こし、その抑制性ニューロンが標的にブレーキをかける。符号の反転には、ニューロン1個と小さな遅延1回分のコストがかかる。

\section*{「ただしBがなければ」}

ブレーキを使った最も役に立つ演算は禁止である。「Aで鳴れ、ただしBが来ていなければ」。ニューロンはAから押され、Bが起こす仲介役からブレーキを受ける。禁止と「OR」があれば、どんな論理でも組み立てられる。

\pencilfig{3}{\pencilart{3}}{「A、ただしBがなければ」：赤いニューロンが青いニューロンの道をふさぐ。}

時間の中での禁止は、あなたの目にもある動き検出器を生む。隣り合う2つのセンサーのうち、一方は出力ニューロンを興奮させ、もう一方は遅れてそれを抑制する。光の点が一方向に適切な速さで動くと、抑制はちょうど興奮と同時に届き、それを打ち消す。逆方向なら抑制は間に合わず、ニューロンはカチッと鳴る。こうして、一方向の動きだけを見るニューロンができる。

\pencilfig{103}{%
% sensors on top; the left one excites the output, the right one inhibits it through a GABA go-between and a delay
\draw[dotpen=pcAmber, wash=pcAmber] (0,1.6) circle (0.16); \draw[dotpen=pcAmber, wash=pcAmber] (1.6,1.6) circle (0.16);
\node[lbl, anchor=south] at (0.8,1.8) {2つのセンサー};
\draw[pen=pcBlue, wash=pcBlue] (0.4,0) circle (0.24);
\draw[arr=pcBlue] (0,1.44) -- (0.3,0.25);
\draw[pen=pcBlue] (1.6,1.44) -- (1.6,1.18);
\draw[penlite, fill=white] (0.95,0.9) rectangle (2.25,1.18); \node[font=\tiny\itshape, text=pcGray] at (1.6,1.04) {遅延};
\draw[pen=pcBlue] (1.6,0.9) -- (1.6,0.72);
\draw[pen=pcRed, wash=pcRed] (1.6,0.55) circle (0.17);
\draw[pcRed, line width=0.7pt, -{Bar[width=5pt]}] (1.45,0.45) -- (0.66,0.08);
\draw[arr] (0.4,-0.24) -- (0.4,-0.6); \node[lbl, anchor=north] at (0.4,-0.6) {出力};
% outcomes
\begin{scope}[xshift=3.1cm]
  \draw[dotpen=pcAmber, wash=pcAmber] (0,1.25) circle (0.1); \draw[arr=pcAmber] (0.18,1.25) -- (1.1,1.25);
  \node[lbl, anchor=west] at (1.25,1.25) {左から右：カチッ};
  \draw[dotpen=pcAmber, wash=pcAmber] (1.1,0.45) circle (0.1); \draw[arr=pcAmber] (0.92,0.45) -- (0,0.45);
  \node[lbl, anchor=west] at (1.25,0.45) {右から左：沈黙};
  \node[lbl, anchor=west, align=left] at (0,-0.35) {抑制が興奮と\\ちょうど同時に届く};
\end{scope}
}{動き検出器：抑制は遅れて届き、点が右から左へ動くときだけ応答を打ち消す。}


\section*{引き算か、割り算か}

抑制は引き算だけではない。開いた抑制性チャネルは電圧を下へ引っぱるのではなく、\emph{静止状態へ}引き戻す。バケツにもう一つ穴をあけたようなものだ。そのため、どれだけ水が流れ込んでも、水位の上がり方は小さくなる。抑制は信号から一定量を差し引くのではなく、信号を\emph{割る}。これは音量つまみである。抑制の強さが周囲全体の活動とともに増えれば、各ニューロンは自分の信号がどれだけ強いかではなく、\emph{隣と比べて}どれだけ強いかに反応する。だから同じネットワークが、薄暗がりでも真夏の日差しの下でも働く。ただし、スパイクの頻度に対してこの「割り算器」がきれいに効くのは、生きた脳に常にある一定の背景ノイズと組み合わさったときだけだ。

\pencilfig{104}{%
\foreach \dx/\t in {0/引き算,4.6/割り算}{
  \begin{scope}[xshift=\dx cm]
  \draw[pen] (0,0) -- (3.6,0); \draw[pen] (0,0) -- (0,1.9);
  \node[lbl, anchor=north] at (1.8,-0.05) {信号}; \node[lbl, anchor=south, rotate=90] at (-0.05,0.95) {応答};
  \draw[pen=pcBlue] (0.4,0) -- (3.3,1.75);
  \node[font=\small\itshape, text=pcGray, anchor=south] at (1.8,1.95) {\t};
  \end{scope}}
\draw[pen=pcRed, line width=0.9pt] (1.3,0) -- (3.4,1.27);
\node[lbl, text=pcRed, anchor=west] at (2.25,0.35) {ずれ};
\begin{scope}[xshift=4.6cm]
\draw[pen=pcRed, line width=0.9pt] (0.4,0) -- (3.3,0.8);
\node[lbl, text=pcRed, anchor=west] at (2.0,0.25) {傾きが小さい};
\end{scope}
\node[lbl, text=pcBlue, anchor=east] at (2.25,1.45) {抑制なし};
}{青は抑制なしのニューロンの応答、赤は抑制ありの応答。引き算は閾値をずらし、割り算は音量つまみのように音量を下げる。}

もう一つの帰結がある。余分な穴があると、バケツは早く忘れる。2つの信号が足し合わされるために入らなければならない窓が狭くなる。抑制は同時検出器をより厳しくするのだ。

\section*{選択}

いよいよ、足りなかった一番大事なもの、多くの中から1つを選ぶことだ。2つのニューロン群が互いに抑制し合うとしよう。抑制が弱いうちは両方が働き、ただ静かになるだけだ。しかし抑制が十分強ければ、どんな小さな差もシーソーのように増幅され、ついには一方の群れが完全に黙る。勝者総取りである。しかも勝利は保たれる。ネットワークが考えを変えるには、新しい根拠が古いものより目に見えて強くなければならない。マシンは物音ひとつで動揺したりしない。

ブレーキは記憶も消す。抑制性ニューロンは「リセット」信号を受けると、前の章のたき火を消す。互いに抑制し合う2つのこうしたセルは、あらゆるコンピュータメモリの基礎であるラッチに相当する。

\section*{ぎりぎりの生活}

「プラス」だけのネットワークは雪崩に陥ることがある。1つのスパイクが次のスパイクを1つより多く生むのだ。ブレーキは、オーブンの温度を保つ制御器のように、ネットワークを動作点に保つ。条件は2つある。抑制は十分に強く、そして十分に速くなければならない。強くても遅ければ、制御器はやりすぎ始める。お湯が届くのを待たずにシャワーの蛇口をひねり続ける人のようなものだ。温度は行ったり来たりする。脳では、「プラス」と「マイナス」のこうしたループが速いリズム、毎秒数十回の振動を生む。

大脳皮質は、どうやらまさにぎりぎりのところで生きている。自身のフィードバックは、ブレーキがなければ雪崩になるほど強い。それを抑えているのは速いブレーキだけであり、だからこそ皮質は弱い信号にもこれほど敏感に応える。

この章の結論は意外なものだ。ブレーキはマシンに新しい計算をほとんど加えない。加えるのは\emph{制御}である。選択、リセット、音量調整、安定性。興奮がデータを運び、抑制はどのデータを、いつ、どれだけの大きさで通すかを決める。

\fullversion{3}
